% Cote 119: offprints, reports, courses, CV and letters (typed and printed part): archive cote 119, site page 200 (PDF page 201).
% typescript. Draft compared: none.
% AI-transcribed by Claude Opus 5.5 (claude-opus-5-5, adaptive thinking, effort high), Message Batches API, batch r119_full (msgbatch_01DwZNm5Y7YsKBWyAN4sFs8w), 2026-11-13; not yet checked by eye.
\scanpage{119/200}
\mspage{- 5 -}

théorème qui s'énonce en termes de catégories dérivées (formulation sur laquelle
j'avais insisté de longue date, et qui a été prouvée indépendamment par R. Kiehl
[*] et O.Forster et K.Knorr [*]). D'autres théorèmes de finitude(de Frisch et Siu)
pour les images directes supérieures d'un faisceau cohérent par une immersion
ouverte, utilisant la profondeur du faisceau en les points du complémentaire,
sont inspirés de théorèmes analogues en géométrie algébrique [16, SGA 2];
remarques analogues pour des théorèmes sur la cohomologie à supports compacts
des faisceaux algébriques cohérents, complétés par un théorème d'existence,
et leur interprétation en termes de théorèmes du type de Lefschetz pour la
cohomologie cohérente (la version algébrique fait partie de la thèse de Mme
Michèle Raynaud (en cours de publication), et la version analytique est due
à Trautmann [*]). Parlant en termes de grands thèmes de recherche plutôt qu'en
termes de résultats techniques particuliers, je pense qu'outre les thèmes déjà
nommés,les thèmes suivants ont été directement suscités ou tout au moins influ-%
encés par les idées que j'avais développées en géométrie algébrique :

a) \uline{Techniquesde construction d'espaces analytiques}, aboutissant aussi bien à
des espaces ``modulaires'' ``globaux'' comme les espaces modulaires de Picard,
pour certains espaces analytiques compacts comme dans [11] (le cas général ne
semble pas encore traité), qu'à des espaces modulaires ``locaux'' de déformation
d'une structure analytique complexe donnée, ou au modèle de la Géométrie For-%
melle (``th. d'existence des modules formels'', cf. [15 bis,Exp. n° 195]).
Dans certains cas, les énoncés obtenus en géométrie algébrique sont directement
applicables (cf. M. Hakim [*]), dans d'autres de nouvelles difficultés surgis-%
sent, pas toujours surmontées à l'heure actuelle. Parmi les travaux définitifs
dans ce sens, on peut citer la thèse de A. Douady [*].

b) \uline{Théorèmes de dualité locaux et globaux pour les faisceaux cohérents}, déve-%
loppés notamment par J.L. Verdier [*] et J.P. Ramis et G. Ruget [*], inspirés
par la théorie que j'avais développée dans le cas des schémas, exposée dans
R. Hartshorne [*].

c) \uline{Formulations de théorèmes du type de Riemann-Roch} pour des variétés analy-%
tiques compactes ou des morphismes propres de telles variétés, cf. [16, SGA 6,
Exp. 0]. Les problèmes essentiels restent toujours ouverts.

d) \uline{Théorèmes de De Rham analytiques complexes} [16 bis], \uline{cohomologie cristalline}
\uline{complexe}. Certains des résultats et des idées que j'avais développés à ce
sujet ont été utilisés dans des développements théoriques divers, comme la
théorie de Hodge généralisée de P. Deligne [*].

\mspage{\hand{200.}}
